\subsection{The unit connector with terminal arms of lengths two and three}\label{proof:R2-C007}
We use the auxiliary coordinates $B=-K$ and $u=B\one=-g$. Thus
$T=J+B=A^*$, $A=J-B$, $p=1-u$, $q=1+u$, and
$\bar f_n=T^n\one$, $h_n=Tf_{n-1}$, $\delta_i=\|B^i\one\|_2^2$,
$d=\delta_1$, $\theta=13/32$, and $\tau=d-\delta_2$.
All expectations and adjoints use the original measure. The two residual
terminal arms have lengths $(2,3)$ and messages $h_2,f_3$; the unit
connector points from the terminal to the receiving root and has output
$X=T(h_2f_3)$. The receiving arm has its actual outward message $f_a$,
where $a\ge3$. These directions depend only on the original arm lengths.
We prove
\begin{equation}\tag{R07.0.1}\label{eq:R07:claim}
 \E[f_aX^4]\le1-d/100.
\end{equation}
The same conclusion holds with $\E[E_of_a\{T(E_th_2f_3)\}^4]$ on the left,
where $E_o,E_t$ are any actual equal-arm decorations.

The finite calculus of Q3-C014 gives
$0\le d\le1/3$, $d^2\le\delta_2\le\theta d$,
$\delta_{i+1}\le\theta\delta_i$, $\|B\|^2\le\theta$,
$|BV(x)|\le\|V\|_2$, $f_n\le2$, and $\E f_n\le1-2d$ for $n\ge3$.
The operators $A,T$ are positive. Equal-arm pairs lie in $[0,1]$,
and $p,q>0$ on the original positive support.
Put $G_n(z)=\sum_{j=0}^{n-1}z^j$, $\beta_1=\beta_2=0$, and
$\beta_n=\sum_{j=1}^{\lfloor(n-1)/2\rfloor}(-1)^{j-1}\delta_j$.
The exact positive convolution \eqref{eq:Q117:conv} and the receiving uppers
\eqref{eq:Q117:receiver} and \eqref{eq:Q187:receiver} give the receiving bounds
\begin{equation}\tag{R07.1.1}\label{eq:R07:receiving}
 \begin{gathered}
 S_a=1-G_a(-B)u,\qquad C_3=f_3=S_3-d,\\
 C_a=S_a-\beta_a-\beta_{a-1}p\quad(a\ge4),\qquad
 0\le f_a\le C_a,\qquad \tau\le\beta_j\le d\quad(j\ge3),\\
 C_a\le(1-d)p+\tfrac65\sqrt{pqd}
       \le(1-d)p+\tfrac65\sqrt d\quad(a\ge3).
 \end{gathered}
\end{equation}
For $a=3$ the last bound uses the actual variance
$d+\delta_2-d^2\le36d/25$ of $u-Bu$. For $a\ge4$ it retains two
convolution terms, uses $\beta_a+\beta_{a-1}\ge2\tau\ge d$, and
$\|G_{a-1}(-B)u\|_2\le\sqrt d$.
On the original spectral measure $\nu$ of $-B^2$ at $u$, write
$z=i\sqrt x$, $0\le x\le\theta$, and
$\int x^j\,d\nu=\delta_{j+1}$. If $G_a=E_a+zD_a$, then
$1-x\le E_a,D_a\le1$ for $a\ge3$.

Put $m=1-d/2$ and $s=-B^2u$. Since $h_2+f_3=2m+s$,
\begin{equation}\tag{R07.2.1}\label{eq:R07:Y}
 0\le X\le T(m+s/2)^2\le Y:=m^2q-mB^3u+\delta_3/2.
\end{equation}
Here $\E s=0$, $Ts=-B^3u$, and $T(s^2)\le2\delta_3$; in particular,
$Y\ge0$. With $w=Y-m^2$ and $c=\delta_3/2$ the exact moments are
\begin{equation}\tag{R07.2.2}\label{eq:R07:moments}
 \E w=-m\delta_2+c,\qquad
 \E w^2=m^4d+m^2\delta_4-m\delta_2\delta_3+\delta_3^2/4.
\end{equation}
The mixed product of $u$ and $B^3u$ vanishes; the negative constant--mean
cross term in \eqref{eq:R07:moments} is retained. We use the identity
\begin{equation}\tag{R07.2.3}\label{eq:R07:fourth}
 Y^4=m^8+4m^6w+w^2(Y^2+2m^2Y+3m^4).
\end{equation}

For $0\le d\le1/4$ the following continuous pointwise bound holds:
\begin{equation}\tag{R07.3.1}\label{eq:R07:curvature}
 C_a(Y^2+2m^2Y+3m^4)\le8m^2\quad(a\ge3).
\end{equation}
Here is its full certificate. Set $t=\sqrt d\in[0,1/2]$,
$\alpha=1-t^2$, $c_0=6t/5$, and $C=5t/12+t^2/10$.
The row bound and $\delta_3\le\theta^2d$ give $Y\le m^2q+mC$:
indeed $\theta<5/12$, $\theta^2/(2m)\le4\theta^2/7<1/10$, and
$m\ge7/8$. After division by $m^2$ it suffices to bound
\[
 Q(q,t)=[\alpha(2-q)+c_0]\{[m(q+1)+C]^2+2m^2\}\le8,
 \qquad0\le q\le2.
\]
Let $h=C/m\le4/15$. The derivative $Q_q$ is a concave quadratic, and
\[
 \frac{Q_q(0)}{m^2}=\alpha(1+2h-h^2)+2c_0(1+h)>0,\quad
 \frac{Q_q(1/2)}{m^2}=\alpha(1/4-h^2)+2c_0(3/2+h)>0.
\]
Thus $Q$ increases on $[0,1/2]$. On $q\ge1/2$,
\[
 Q_{qq}=2m^2(c_0-3\alpha q)-4\alpha mC\le-5/4.
\]
For $t\le2/5$ the first term is at most $-41262/31250<-5/4$.
For $t\ge2/5$, use $m^2\ge49/64$, $c_0-3\alpha/2\le-21/40$,
$\alpha m\ge21/32$, and $C\ge137/750$ to obtain
$-Q_{qq}\ge1029/1280+2877/6000>5/4$.
The same formula gives $Q_{qq}\le-\lambda$ with
\[
\begin{array}{c|ccccc}
 t\text{ upper bound}&1/8&1/5&1/4&1/3&1/2\\\hline
 \lambda&5/2&9/4&2&3/2&5/4
\end{array}
\]
by inserting the respective bounds for $c_0,\alpha,m$ and retaining the
nonpositive final term. Direct multiplication gives
\begin{align*}
 Q(1,t)&=6+\tfrac{133}{15}t-\tfrac{6787}{720}t^2
 -\tfrac{5357}{600}t^3+\tfrac{21011}{3600}t^4
 +\tfrac{1161}{500}t^5-\tfrac{131}{100}t^6\\
 &\le F(t):=6+\tfrac{133}{15}t-\tfrac{159}{20}t^2-\tfrac{25}{3}t^3,\\
 D(t):=Q_q(1,t)&=-2+\tfrac{119}{30}t+\tfrac{3331}{720}t^2
 -\tfrac{509}{150}t^3-\tfrac{9131}{3600}t^4
 +\tfrac{56}{75}t^5+\tfrac{41}{100}t^6.
\end{align*}
The first comparison uses $t^4\le t^2/4$, $t^5\le t^3/4$.
$F$ increases on $[0,1/3]$, $F'(1/3)=71/90$, and
$F''\le-159/10$, so
$F\le12577/1620+(71/90)^2(5/159)<39/5$ on the full interval.
Discarding positive derivative terms and bounding the two negative
ones at $t=1/2$ gives $D'\ge1105/7200>0$.
Exact endpoint substitution gives
$D(0)=-2$, $D(1/8)>-3/2$, $D(1/5)>-16/15$,
$D(1/4)>-4/5$, $D(1/3)>-1/3$, and $D(1/2)<3/5$.
The strong concavity bound $Q\le Q(1,t)+D(t)^2/(2\lambda)$ therefore
has the following certificates on five closed intervals:
\[
\begin{array}{c|c|c|l}
 t&|D|&\lambda&\text{upper bound, strictly below }8\\\hline
 {}[0,1/8]&2&5/2&F(1/8)+4/5\\
 {}[1/8,1/5]&3/2&9/4&F(1/5)+1/2\\
 {}[1/5,1/4]&16/15&2&F(1/4)+64/225\\
 {}[1/4,1/3]&4/5&3/2&F(1/3)+16/75\\
 {}[1/3,1/2]&3/5&5/4&39/5+18/125
\end{array}
\]
These rational comparisons, including $12577/1620+16/75<8$ in the
fourth row, prove \eqref{eq:R07:curvature} throughout its domain.

For $a=3$ and $0<d\le1/4$, put $v=\delta_2/d$,
$w_3=\delta_3/d$, $w_4=\delta_4/d$. Then
$d\le v\le\theta$, $v^2\le w_3\le\theta v$, $w_4\le\theta w_3$,
and the actual receiving moments are
\begin{align*}
 \E C_3&=1-2d,\\
 \E[C_3(m^2u-mB^3u)]&=-m^2(d-\delta_2)-m(1-d)\delta_2+m\delta_3,\\
 \E[C_3w]/d&=-m^2+(md/2)v+\tfrac32(1-d)w_3.
\end{align*}
Equations \eqref{eq:R07:moments}--\eqref{eq:R07:curvature} bound
$\E[C_3Y^4]$ by
\[
 m^8(1-6d)+8m^6d+2m^7vd^2+6m^6(1-d)w_3d
 +8m^4w_4d-8m^3vw_3d^2+2m^2w_3^2d^2.
\]
Writing its normalized deficit as $D_3=(1-\text{bound})/(m^4d)$,
the positive Taylor expansion gives
\[
 \frac{m^{-4}-m^4(1-6d)}d-8m^2
 \ge2-3d+10d^2-\tfrac78d^3.
\]
Use the spectral bounds above and $m^2\ge49/64$, and discard the
positive term $8vw_3d/m$. This gives
\begin{align*}
 D_3&\ge2-3d+10d^2-\tfrac78d^3-2\theta d
 -[6m^2(1-d)+8\theta]\theta^2-\tfrac{128}{49}\theta^4d\\
 &\ge\tfrac{15}{32}-2d+\tfrac{17}{2}d^2
 \ge\tfrac{191}{544}>1/3.
\end{align*}
For the intermediate comparison the constant is
$2-6\theta^2-8\theta^3=1939/4096>15/32$,
the linear coefficient is greater than $-2$, and, after absorbing the
negative cubic using $d\le1/4$, the quadratic coefficient is
$313/32-(57/8)\theta^2>17/2$. Completing the square proves the last
inequality. Hence the actual deficit exceeds $m^4d/3>d/6$.

For $a\ge4$ work on $0<d\le1/5$, set $\beta=\beta_{a-1}$,
$e_a=\int E_a\,d\nu$, $l_a=\int x^2D_a\,d\nu$. The exact finite
receiving data are
\begin{equation}\tag{R07.5.1}\label{eq:R07:covariance}
\begin{gathered}
 \E C_a=1-\beta_a-\beta_{a+1}-\beta,\qquad
 \E[C_au]=-e_a+\beta d,\\
 \E[C_aB^3u]=(1-\beta_a-\beta)\delta_2-l_a,\qquad
 \tau\le e_a\le d,\quad0\le l_a\le\delta_3.
\end{gathered}
\end{equation}
In $m^8\E C_a+4m^6\E[C_aw]$ the coefficients of $\beta_a$ and
$\beta_{a+1}$ are negative: for the former use
$m^2+2\delta_3-4m\delta_2\ge m(m-4\theta d)>0$.
The coefficient of $\beta$ is $m^6J$, where
$J=-m^2(1-4d)-2\delta_3+4m\delta_2$.
Lowering all three betas and $e_a$ to $\tau$ and raising $l_a$ to
$\delta_3$ costs at most $m^6\delta_2J_+$ beyond the resulting bound $M_0$.
Its normalized deficit satisfies
\begin{align*}
 D_0:=\frac{1-M_0}{m^4d}\ge{}&B_0(d)-3m^4v-2m^3vd(1+d-4v)
 -[m^2(6-8d)+8\theta]w_3\\
 &+vw_3d(8/m-6m^2)-2w_3^2d/m^2,\\
 B_0(d):={}&\frac{m^{-4}-m^4(1-7d+4d^2)}d-8m^2\\
 \ge{}&3-9d+\tfrac{39}{2}d^2-\tfrac{59}{8}d^3+\tfrac{67}{16}d^4.
\end{align*}
The right side decreases in $w_3$: its negative main coefficient is
at least $8\theta$, whereas its positive correction is at most
$vd(8/m)\le13/18$. Insert $w_3=\theta v$. The resulting expression is
$B_0-Lv+dQv^2$, with
\[
 L=3m^4+2m^3d(1+d)+\theta[m^2(6-8d)+8\theta],\quad
 Q=8m^3+\theta(8/m-6m^2)-2\theta^2/m^2<12.
\]
Its $v$ derivative is at most $-3m^4+24d\theta<0$, because
$3m^4\ge3(9/10)^4>39/20\ge24d\theta$. Thus $v=\theta$ gives
\begin{align*}
 D_0\ge{}&B_0-3\theta m^4-2\theta m^3d(1+d-4\theta)
 -\theta^2[m^2(6-8d)+8\theta]\\
 &+\theta^3d(8/m-6m^2)-2\theta^4d/m^2\\
 \ge{}&P(d):=\tfrac{1043}{4096}-\tfrac{11}{3}d
                +\tfrac{151}{10}d^2-\tfrac{49}{10}d^3.
\end{align*}
Indeed $m^{-2}\le1+6d/5$ and
$8/m-6m^2\ge2+10d+d^2/2+d^3$; terms of degree at least four are
nonnegative, while the first four coefficients are bounded by
\[
\begin{gathered}
c_0=1043/4096,\quad
c_1=-9+\tfrac{29}{4}\theta+14\theta^2+2\theta^3-2\theta^4\ge-11/3,\\
c_2=\tfrac{39}{2}-\tfrac{67}{8}\theta-\tfrac{19}{2}\theta^2
       +10\theta^3-\tfrac{12}{5}\theta^4\ge151/10,\\
c_3=-\tfrac{59}{8}+\tfrac{87}{16}\theta+2\theta^2+\theta^3/2\ge-49/10.
\end{gathered}
\]
For example $c_1+11/3=3373/1572864$ and
$c_2-151/10=9145/262144$.
The positive-beta cost is controlled using $\delta_3\ge\delta_2^2/d$:
\[
 J\le-m^2(1-4d)+(4mv-2v^2)d
 \le J_*(d):=-1+\tfrac{3223}{512}d-\tfrac{81}{16}d^2+d^3.
\]
The second comparison uses $m>v$. The polynomial $J_*$ increases on
$[0,1/5]$, $J_*(9/50)<0$, and $J_*(1/5)=4127/64000<1/15$.
For $d\le9/50$ the cost vanishes, and
\[
 P(d)\ge1043/4096-(11/3)d+(7109/500)d^2>1/64.
\]
Completing the square reduces this to
$1043/4096-15125/63981>1/64$, whose cross-multiplied positive
difference is $685399$. On $[9/50,1/5]$, $P$ increases and
$P(9/50)>1/20$, whereas the normalized extra cost is
$m^2vJ_+<\theta/15<1/36$. The remaining deficit is at least
$1/20-1/36=1/45>1/64$. Consequently
\[
 \E[C_aY^4]\le1-(m^4/64)d\le1-d/100.
\]

For the remaining range use the actual terminal, whose exact mean is
$\E[h_2f_3]=1-2d-d^2$. This follows by multiplication, or by subtracting
the means of the squares in
\[
h_2f_3=(m+s/2)^2-(-u+Bu-B^2u/2-d/2)^2;
\]
the two square means are $m^2+\delta_3/4$ and
$d+\delta_3/4+5d^2/4$. Positivity gives $X\le2(1-2d-d^2)$.
For $1/5\le d\le1/3$ both nonnegative factors decrease, so
\[
 \E[f_aX^4]\le(1-2d)[2(1-2d-d^2)]^4
 \le\tfrac35(28/25)^4=1843968/1953125<19/20\le1-d/100.
\]
This overlaps both low ranges and proves \eqref{eq:R07:claim}.
Decorations follow from $0\le E_o,E_t\le1$ and positivity of $T$.
At $d=0$, $u=0$ and successive leaf integration gives every original
tree message equal to one; the kernel need not vanish. All original
arms, decorations, and the unit connector retain their original edges.
Full reversal changes the receiving arm, terminal arms, connector, and
decorations together. Combining this unit row with \eqref{eq:q117-unified:shorttwo}
at $L\ge2$ gives, for terminal $(2,3)$, every $L\ge1$ with deficit
$d/400$. Terminals $(2,t)$ with $t\ge4$ require a separate estimate.
